\documentclass[11pt]{article}
\usepackage[a4paper,margin=1in]{geometry}
\usepackage[T1]{fontenc}
\usepackage{lmodern}
\usepackage{amsmath,amssymb,amsthm,mathtools}
\usepackage{microtype}
\usepackage[hidelinks]{hyperref}
\newtheorem{theorem}{Theorem}[section]
\newtheorem{lemma}[theorem]{Lemma}
\newtheorem{definition}[theorem]{Definition}
\theoremstyle{remark}
\newtheorem{remark}[theorem]{Remark}
\newcommand{\Xt}{\widetilde X}
\newcommand{\Lt}{\widetilde L}
\newcommand{\wt}{\widetilde\omega}
\newcommand{\OO}{\mathcal O}
\newcommand{\Q}{\mathcal Q}
\newcommand{\PP}{\mathcal P}
\newcommand{\Ioc}{I_{\mathrm{oc}}}
\newcommand{\Rcub}{\mathcal R^{\mathrm{cub}}}
\newcommand{\EE}{\mathcal E}
\DeclareMathOperator{\ord}{ord}
\DeclareMathOperator{\dv}{div}
\DeclareMathOperator{\Per}{Per}
\DeclareMathOperator{\Supp}{Supp}
\title{Spin Parity and Period Obstructions}
\author{Chuyang Chen}
\date{October 2026}
\newcommand{\publicationstatus}{%
\begin{center}
\fbox{\begin{minipage}{0.92\linewidth}
\small
\textbf{AI EXPLORATORY MATHEMATICAL NOTE}\\[2pt]
\textbf{Status:} AI-assisted; \textbf{NO INDEPENDENT HUMAN REVIEW}\\
\textbf{Version:} v1.1 (publication-format revision)\\
\textbf{First repository publication:} October 24, 2026\\
\textbf{Related Lean formalization:} Conditional F1. Internal arguments are checked relative to explicit external results and geometric interfaces; see the end matter for the exact scope.
\end{minipage}}
\end{center}
\par\medskip
}
\hypersetup{pdftitle={Spin Parity and Period Obstructions},pdfauthor={Chuyang Chen}}
\begin{document}
\maketitle
\publicationstatus
\begin{abstract}
This note addresses the two parts of Chen--Yu's Problem 3.2 for quadratic differentials.
It gives a section-count interpretation of spin parity on the canonical double cover, and it
gives a flat orientation-cover invariant distinguishing the regular and irregular exceptional
components, with the component identification supplied by Chen--M\"oller's theorem.
\end{abstract}
\begin{quote}
\small
\textbf{Feedback.} This note has not undergone independent human mathematical review.
Readers who know relevant prior literature, identify an error, or wish to check the argument
are especially encouraged to contact the author listed at the end of the paper or to open an issue in the repository.
\end{quote}

\section{Introduction}
Problem 3.2 in the survey of Chen and Yu \cite{ChenYu} asks for two refinements in the topology of strata
of quadratic differentials. The first is to recover the spin parity of the canonical double cover
by computing the dimension of the corresponding theta-characteristic. The second is to find
a flat-geometric invariant distinguishing the regular and irregular connected components of the
exceptional strata in genera three and four.

The exceptional strata under discussion are
\[
\Q(9,-1),\quad \Q(6,3,-1),\quad \Q(3,3,3,-1)
\]
in genus three, and
\[
\Q(12),\quad \Q(9,3),\quad \Q(6,6),\quad \Q(6,3,3),\quad \Q(3,3,3,3)
\]
in genus four. Chen and M\"oller proved that the genus-three exceptional strata have two components,
called regular and irregular, distinguished by a Brill--Noether condition involving one
third of the zero divisor \cite[Theorem 6.1]{ChenMoller}. In genus four, \cite[Theorem 7.1]{ChenMoller} gives the analogous
distinction on the two non-hyperelliptic regular/irregular components. The strata $\Q(6,6)$,
$\Q(6,3,3)$, and $\Q(3,3,3,3)$ also have hyperelliptic components; statements below about regular
and irregular genus-four components are made on the non-hyperelliptic part.

The second part of this note uses Chen--M\"oller's component theorem as an input. The new
point is to realize their Brill--Noether condition as an invariant on the canonical orientation
double cover and as an anti-invariant Mittag--Leffler period obstruction. Thus the result is a
conservative solution of the flat-geometric part of Problem 3.2, not an independent reproof of
the Chen--M\"oller component theorem.

\section{Spin parity from theta-characteristic sections}
Let $(X,q)\in\Q(m_1,\ldots,m_n)$ be a primitive quadratic differential with at most simple poles, so
$m_i\geq -1$. Assume that
\[
m_i\not\equiv 2\pmod 4
\]
for all $i$. Let
\[
\pi:\Xt\longrightarrow X
\]
be the canonical double cover and write
\[
\pi^*q=\wt^2.
\]
The cover is ramified exactly over odd-order singularities, including simple poles. Since $q$ is
primitive, the cover is connected. Under the congruence assumption all singularities of $\wt$ have
even order, so
\[
\Lt=\OO_{\Xt}\left(\frac12\dv(\wt)\right)
\]
is a theta-characteristic. We use the standard convention that the parity of a theta-characteristic
$L$ is $h^0(L)\bmod 2$. Lanneau's constancy result and parity formula are stated via Arf-invariant
methods \cite{Lanneau}; the following proof gives the requested $h^0$-interpretation.

\begin{theorem}[Spin parity formula]\label{thm:spin}
With the notation above,
\[
h^0(\Xt,\Lt)\equiv 1-g+\sum_i\left\lfloor\frac{m_i+1}{4}\right\rfloor\pmod 2.
\]
\end{theorem}
\begin{proof}
Write
\[
\dv(q)=\sum_i m_i z_i.
\]
Set
\[
A=\sum_i\left\lfloor\frac{m_i}{2}\right\rfloor z_i,\qquad
B=\sum_{m_i\ \mathrm{odd}}z_i.
\]
Then
\[
\dv(q)=2A+B\sim 2K_X.
\]
Let
\[
M=K_X(-A).
\]
Then
\[
M^{\otimes 2}\cong\OO_X(B),
\]
and the canonical double cover satisfies
\[
\pi_*\OO_{\Xt}\cong\OO_X\oplus M^{-1}.
\]
Define
\[
D_+=\sum_i\left\lfloor\frac{m_i+1}{4}\right\rfloor z_i.
\]
We prove
\begin{equation}\label{eq:splitting}
\pi_*\Lt\cong\OO_X(D_+)\oplus K_X(-D_+).
\end{equation}
The divisor $\frac12\dv(\wt)$ is invariant under the deck involution $\tau$, so $\Lt$ carries its natural $\tau$-linearization.
We decompose $\pi_*\Lt$ into invariant and anti-invariant summands.

The calculation is local. If $m=4r$, the cover is unbranched and locally $\wt=\pm z^{2r}dz$ on each
sheet. Thus $\Lt$ has coefficient $r$ above the point. The invariant part gives $\OO_X(rz)$. Multiplication
of an anti-invariant local section by $\wt$ descends to a differential vanishing to order $r$, giving
$K_X(-rz)$. This is the desired contribution because $\lfloor(4r+1)/4\rfloor=r$.

If $m=4r+1$, put $z=w^2$. Then
\[
\pi^*q=(w^2)^{4r+1}(2w\,dw)^2=4w^{8r+4}dw^2,
\]
so $\ord(\wt)=4r+2$, and $\Lt$ has coefficient $2r+1$ at the ramification point. An invariant local section
has only even Laurent terms, and the largest possible pole is $w^{-2r}=z^{-r}$, yielding $\OO_X(rz)$. The
largest anti-invariant pole is $w^{-(2r+1)}$. Multiplying by $\wt=2w^{4r+2}dw$ gives $2w^{2r+1}dw=z^r dz$, so
the anti-invariant summand gives $K_X(-rz)$. This matches $D_+=rz$.

If $m=4r+3$, again $z=w^2$ and
\[
\pi^*q=(w^2)^{4r+3}(2w\,dw)^2=4w^{8r+8}dw^2.
\]
Thus $\ord(\wt)=4r+4$, and $\Lt$ has coefficient $2r+2$. The invariant part permits even Laurent
poles up to order $2r+2$, hence gives $\OO_X((r+1)z)$. The highest allowed Laurent pole is even
and therefore does not occur in the anti-invariant summand. The largest anti-invariant pole is
$w^{-(2r+1)}$, and
\[
w^{-(2r+1)}\cdot 2w^{4r+4}dw=2w^{2r+3}dw=z^{r+1}dz.
\]
Thus the anti-invariant summand is $K_X(-(r+1)z)$, matching $D_+=(r+1)z$.

If $m=-1$, locally $q=z^{-1}dz^2$. With $z=w^2$, one has $\pi^*q=4dw^2$, so $\wt$ is regular and nonvanishing.
Hence $D_+$ receives no contribution. The invariant part gives $\OO_X$, and an anti-invariant
holomorphic function is $wh(w^2)$. Multiplying by $dw$ gives $\frac12 h(z)dz$, an arbitrary holomorphic
differential on the base. Thus the anti-invariant summand is $K_X$. Finally, if $m=0$, the cover is
unbranched and $\wt$ is nonvanishing; the local splitting is $\OO_X\oplus K_X$.

These local computations prove \eqref{eq:splitting}. Therefore
\[
h^0(\Xt,\Lt)=h^0(X,\OO_X(D_+))+h^0(X,K_X(-D_+)).
\]
By Riemann--Roch,
\[
h^0(X,\OO_X(D_+))-h^0(X,K_X(-D_+))=\deg D_++1-g.
\]
Modulo $2$, sum and difference agree. Hence
\[
h^0(\Xt,\Lt)\equiv\deg D_++1-g
=1-g+\sum_i\left\lfloor\frac{m_i+1}{4}\right\rfloor\pmod 2.
\]
\end{proof}

\section{The orientation-cover invariant}
Let $(X,q)$ be in one of the genus-three exceptional strata, or in one of the non-hyperelliptic
regular/irregular genus-four exceptional strata listed above. Let $Z(q)$ be the zero divisor of $q$,
and define
\[
D_q=\frac13 Z(q).
\]
In genus three, the simple pole is not included in $D_q$. Let
\[
\pi:\Xt\longrightarrow X,\qquad \pi^*q=\omega^2
\]
be the canonical orientation double cover, with deck involution $\tau$. Then $\tau^*\omega=-\omega$. Define
\[
E_q=\frac12\pi^*D_q,\qquad \Ioc(q)=h^0(\Xt,\OO_{\Xt}(E_q))^-.
\]

\begin{lemma}\label{lem:integral}
The divisor $E_q$ is integral.
\end{lemma}
\begin{proof}
If $q$ has zero order $3a$ at $z$, then $D_q$ has coefficient $a$ at $z$. If $a$ is even, then the orientation
cover is unbranched over $z$, and $\frac12\pi^*(az)$ has coefficient $a/2$ at each of the two points above $z$.
If $a$ is odd, the cover is ramified and $\pi^*(az)=2a\widetilde z$, so $\frac12\pi^*(az)=a\widetilde z$. Hence all coefficients are
integral.
\end{proof}

\begin{theorem}[Orientation-cover invariant]\label{thm:orientation}
There is a natural isomorphism
\[
H^0(X,K_X(-D_q))\cong H^0(\Xt,\OO_{\Xt}(E_q))^-.
\]
Consequently,
\[
\Ioc(q)=h^0(X,K_X(-D_q)).
\]
\end{theorem}
\begin{proof}
Define
\[
\Phi:H^0(X,K_X(-D_q))\longrightarrow H^0(\Xt,\OO(E_q))^-
\]
by
\[
\Phi(\eta)=\frac{\pi^*\eta}{\omega}.
\]
It is anti-invariant because $\pi^*\eta$ is invariant and $\omega$ is anti-invariant. We check the pole bound.

If $q$ has zero order $3a$ at $z$ and $a$ is even, the cover is unbranched and $\ord(\omega)=3a/2$ above $z$.
Since $\eta$ vanishes to order at least $a$, the quotient has order at least $-a/2$, exactly the allowed
pole order given by $E_q$. If $a$ is odd, write $z=w^2$. Then
\[
\pi^*q=4w^{6a+2}dw^2,\qquad \omega=2w^{3a+1}dw.
\]
If $\eta=z^s dz$ with $s\geq a$, then $\pi^*\eta=2w^{2s+1}dw$. Therefore
\[
\ord(\pi^*\eta/\omega)\geq (2a+1)-(3a+1)=-a,
\]
again exactly the pole bound prescribed by $E_q$. At a simple pole, $q=z^{-1}dz^2$, and with $z=w^2$
one has $\pi^*q=4dw^2$. Thus $\omega$ is regular and nonvanishing. The pole is not in $D_q$, so $E_q$ has
coefficient zero there, and the quotient is holomorphic.

Conversely, let $f\in H^0(\Xt,\OO(E_q))^-$. Then $f\omega$ is $\tau$-invariant and hence descends to a meromorphic
differential $\eta$ on $X$. The preceding calculations in reverse show that $\eta$ is holomorphic and
vanishes along $D_q$. Indeed, in the unbranched case the order of $f\omega$ is at least $-a/2+3a/2=a$.
In the ramified case the order is at least $-a+(3a+1)=2a+1$, which descends to a differential
$z^s dz$ with $s\geq a$. At a simple pole, since $E_q$ has coefficient zero and $f$ is anti-invariant, locally
$f(w)=c_1w+c_3w^3+\cdots$. Thus
\[
f\omega=2(c_1w+c_3w^3+\cdots)dw=(c_1+c_3z+\cdots)dz,
\]
which descends holomorphically. Hence $\eta\in H^0(X,K_X(-D_q))$, and the two constructions are
inverse.
\end{proof}

\section{Anti-invariant period obstruction}
The invariant in Theorem \ref{thm:orientation} can be formulated as a finite-dimensional Mittag--Leffler obstruction
on the orientation cover. Let $\PP_q$ be the vector space of principal parts of meromorphic functions
with poles bounded by $E_q$ along $\Supp(E_q)$. It is naturally a representation of $\tau$. Let
\[
\delta_q:\PP_q\longrightarrow H^1(\Xt,\OO_{\Xt})
\]
be the Mittag--Leffler coboundary map, and let
\[
\delta_q^-:\PP_q^-\longrightarrow H^1(\Xt,\OO_{\Xt})^-
\]
be its anti-invariant part. Then
\begin{equation}\label{eq:kernel}
\ker\delta_q^-\cong H^0(\Xt,\OO_{\Xt}(E_q))^-.
\end{equation}
\begin{theorem}[Anti-invariant obstruction criterion]\label{thm:obstruction}
In the range of Theorem \ref{thm:orientation},
\[
\ker\delta_q^-\cong H^0(X,K_X(-D_q)).
\]
In particular,
\[
\Ioc(q)>0\quad\Longleftrightarrow\quad\delta_q^-\text{ is not injective}.
\]
\end{theorem}
\begin{proof}
By \eqref{eq:kernel} and Theorem \ref{thm:orientation},
\[
\ker\delta_q^-\cong H^0(\Xt,\OO(E_q))^-\cong H^0(X,K_X(-D_q)).
\]
This proves the theorem.
\end{proof}

There is an equivalent period formulation. Let $P\in\PP_q^-$. Locally $P$ is a principal part of a
meromorphic function, hence $dP$ has zero residues at all points of $\Supp(E_q)$. Choose an anti-invariant
meromorphic differential $\beta_P$ on $\Xt$ whose principal part is $dP$. This is possible by the
usual Mittag--Leffler theorem for differentials of the second kind, because the prescribed residues
are zero. The ambiguity in such a choice is an anti-invariant holomorphic differential.

Fix an anti-invariant symplectic basis
\[
\gamma_1,\ldots,\gamma_{2g^-}
\]
of $H_1(\Xt,\mathbb Z)^-$. The periods of $\beta_P$, modulo the anti-invariant holomorphic ambiguity, define a
linear obstruction map
\[
\Per_q^-:\PP_q^-\longrightarrow
\mathbb C^{2g^-}/\Per(H^0(\Xt,K_{\Xt})^-).
\]
The condition $\Per_q^-(P)=0$ means that one can choose $\beta_P$ with all anti-invariant periods equal to
zero. In that case integration gives a global anti-invariant meromorphic function with principal
part $P$, so $P\in\ker\delta_q^-$. Conversely, if $P$ is the principal part of a global anti-invariant meromorphic
function, then its differential has vanishing periods. Therefore
\[
\ker\Per_q^-=\ker\delta_q^-.
\]
This is the flat period-obstruction form of $\Ioc$.

\section{Residual cubic and the component distinction}
Let $\Q(\mu)$ be one of the exceptional quadratic strata in genus $3$ or $4$. Thus
\[
\mu=(3a_1,\ldots,3a_n,-1),\qquad \sum_i a_i=3,
\]
in genus $3$, and
\[
\mu=(3a_1,\ldots,3a_n),\qquad \sum_i a_i=4,
\]
in genus $4$. In genus $4$ we work on the non-hyperelliptic regular/irregular components. For
$q\in\Q(\mu)$, put
\[
D_q=\frac13 Z(q).
\]
Let $P_q$ be the simple-pole divisor in genus $3$, and put $P_q=0$ in genus $4$. Then
\[
\dv(q)=3D_q-P_q.
\]
\begin{definition}[Residual cubic and positive loci]\label{def:residual}
We say that $q$ admits a residual cubic Abelianization
if there are Abelian differentials $\eta,\theta\in H^0(X,K_X)$ and an effective divisor $R\in X^{(g-2)}$
such that
\[
\dv(\eta)=D_q+R,\qquad \dv(\theta)=3R+P_q,
\]
and
\[
q=\frac{\eta^3}{\theta}.
\]
We denote the residual cubic Abelianization locus by
\[
\Rcub_\mu=\{q:q\text{ admits a residual cubic Abelianization}\},
\]
and the positive orientation-cover locus by
\[
\EE_\mu=\{q:\Ioc(q)>0\}.
\]
\end{definition}

\begin{lemma}[Residual cubic equivalence]\label{lem:residual}
For the exceptional strata above,
\[
h^0(X,K_X(-D_q))>0
\]
if and only if $q$ admits a residual cubic Abelianization. Equivalently,
\[
\EE_\mu=\Rcub_\mu.
\]
\end{lemma}
\begin{proof}
Suppose first that $0\ne\eta\in H^0(X,K_X(-D_q))$. Then
\[
\dv(\eta)=D_q+R
\]
for an effective divisor $R$ of degree $g-2$. Define
\[
\theta=\frac{\eta^3}{q}.
\]
Since $\dv(q)=3D_q-P_q$, we have
\[
\dv(\theta)=3(D_q+R)-(3D_q-P_q)=3R+P_q.
\]
Thus $\theta$ is holomorphic and $q=\eta^3/\theta$.

Conversely, if $q=\eta^3/\theta$ with $\theta$ holomorphic, then
\[
\dv(q)=3\dv(\eta)-\dv(\theta).
\]
At every zero of $q$ belonging to $D_q$, the effectivity of $\dv(\theta)$ forces $\eta$ to vanish at least to the
order prescribed by $D_q$. Hence
\[
0\ne\eta\in H^0(X,K_X(-D_q)).
\]
The equality $\EE_\mu=\Rcub_\mu$ follows from Theorem \ref{thm:orientation}.
\end{proof}

\begin{theorem}[Component distinction]\label{thm:component}
For each exceptional genus-three stratum and for each
non-hyperelliptic regular/irregular genus-four exceptional stratum,
\[
\EE_\mu=\Rcub_\mu=\{q:h^0(X,K_X(-D_q))>0\}.
\]
Equivalently, since $\deg D_q=g$,
\[
q\in\EE_\mu\quad\Longleftrightarrow\quad h^0(X,\OO_X(D_q))\geq 2.
\]
By Chen--M\"oller's theorem \cite[Theorems 6.1 and 7.1]{ChenMoller}, this condition is exactly the traditional
irregular component condition in genus three and on the non-hyperelliptic regular/irregular part
in genus four. Its complement is the corresponding regular component.
\end{theorem}
\begin{proof}
The equality
\[
\EE_\mu=\{q:h^0(X,K_X(-D_q))>0\}
\]
is Theorem \ref{thm:orientation}. The equality with $\Rcub_\mu$ is Lemma \ref{lem:residual}. For the listed exceptional signatures,
$\deg D_q=g$, so Riemann--Roch gives
\[
h^0(X,\OO_X(D_q))-h^0(X,K_X(-D_q))=1.
\]
Thus $h^0(X,K_X(-D_q))>0$ is equivalent to $h^0(X,\OO_X(D_q))\geq 2$. Chen--M\"oller prove that this
Brill--Noether condition distinguishes the exceptional regular and irregular components: in genus
three this is \cite[Theorem 6.1]{ChenMoller}, and in genus four, on the non-hyperelliptic regular/irregular
components, this is \cite[Theorem 7.1]{ChenMoller}. Combining their component theorem with the orientation-cover
and period-obstruction identifications above proves the claim.
\end{proof}

\begin{remark}
The component identification in Theorem \ref{thm:component} uses Chen--M\"oller's Brill--Noether
component theorem as an input. The point of the preceding sections is that the same condition
is naturally visible on the orientation double cover and through anti-invariant flat periods.
\end{remark}

\section{Conclusion}
Theorem \ref{thm:spin} gives the requested section-count interpretation of spin parity. Theorems \ref{thm:orientation}, \ref{thm:obstruction},
and \ref{thm:component} give the flat-geometric invariant for the exceptional regular/irregular distinction: the invariant
is $\Ioc$, equivalently the kernel of the anti-invariant period-obstruction map on the orientation
cover. The final identification with the traditional component names relies on Chen--M\"oller's
component theorem.

\section*{Acknowledgments}
Chuyang Chen is grateful to his advisor, Fei Yu, for his guidance on this project.

\section*{Independent review and Lean verification notice}
\begingroup
\small\sloppy
\begin{center}
\begin{minipage}{0.94\linewidth}
\centering
\textbf{\large No independent human mathematical review is recorded.}
\end{minipage}
\end{center}
\par\medskip
\noindent The accompanying multi-file Lean project compiles under Lean 4.33.1 with Mathlib revision\\
\texttt{0df444a360eaa60ab8c11dca51a86af692955474}.
\par\medskip
\noindent Its status is \textbf{Conditional F1}: internal arguments are formalized relative to explicitly listed
standard external results and geometric representation interfaces. Concrete compact Riemann surface instances satisfying all those interfaces have not been constructed in the development.
\par\medskip
\noindent\textbf{Successful Lean verification establishes correctness of the formalized statement relative to its assumptions. It does not by itself establish exact correspondence between all natural-language claims and their Lean statements.}
\endgroup

\section*{Provenance, AI use, and contact}
\begingroup\small\sloppy
\begin{description}
\item[Author and project coordinator.] Chuyang Chen. This is his independent project.
\item[Contact.] \href{mailto:22535018@zju.edu.cn}{22535018@zju.edu.cn}.
\item[Mathematical manuscript.] The author reports using GPT-5.5 for the original manuscript; a complete prompt log and detailed task allocation are not available.
\item[Lean code generation.] The multi-file project was developed with AI assistance in Codex, including proof revision, external-input organization, builds, and axiom inspection. Earlier sessions' model versions are not established.
\item[Build and publication record.] Checks ran in the assisted workflow; Chuyang Chen ran the publication commands. The full build and entry-point audit passed on October 24, 2026. Initial publication commit: \texttt{88c4a1b}.
\item[Editorial revision.] The supplied manuscript (June 16, 2026) was transcribed into LaTeX using the website's template, with sole-author credit, AI disclosures, and acknowledgments. This is a formatting revision; Lean sources are unchanged.
\item[Independent human mathematical verification.] None recorded.
\item[Independent statement-correspondence audit.] None recorded; no F2 status is claimed.
\end{description}
\noindent This summary is not a full AI transcript. Corrections and independent checks are welcome via the contact above or the repository issue tracker.
\endgroup

\section*{Lean formal verification record}
\begingroup\small\sloppy
\begin{description}
\item[Formalization status.] Conditional F1.
\item[Lean source.] \texttt{notes/chen-yu-problem-3-2/Main.lean}; 32 Lean files in the project.
\item[Principal formal theorem.] \texttt{Problem32.PaperSynthesis.complete\_component\_criterion}.
\item[Build commands.] \texttt{lake build}; \texttt{lake env lean Main.lean}.
\item[Build/kernel result.] PASS in the recorded environment.
\item[Lean version.] Lean 4.33.1.
\item[Library snapshot.] Mathlib revision \texttt{0df444a360eaa60ab8c11dca51a86af692955474}.
\item[Kernel axiom audit.] Final entry-point reports contain only \texttt{propext}, \texttt{Classical.choice}, and \texttt{Quot.sound}, with no \texttt{sorryAx}.
\item[Coverage.] Theorems 2.1, 3.2, 4.1, and 5.3, and Lemma 5.2; the theorem map and detailed coverage are in \texttt{STATUS.md}.
\item[External inputs and scope.] See \texttt{EXTERNAL\_INPUTS.md}. Geometric interfaces remain inputs; genus four is restricted to the non-hyperelliptic regular/irregular part. Axiom reports do not list theorem parameters.
\item[Reproduction.] See \texttt{Verify.ps1}, \texttt{VALIDATION.md}, \texttt{verification.log}, and the source hashes and logs under \texttt{logs/current/}.
\end{description}
\noindent\textbf{Interpretation of the label.} Conditional F1 records kernel acceptance under the documented assumptions; it does not certify concrete geometric instances or an independent human correspondence audit.
\endgroup

\begingroup\raggedright
\begin{thebibliography}{9}
\bibitem{ChenMoller}
D.~Chen and M.~M\"oller,
Quadratic differentials in low genus: exceptional and non-varying strata,
\emph{Ann. Sci. \'Ec. Norm. Sup\'er.} (4) \textbf{47} (2014), no.~2, 309--369;
arXiv:1204.1707.
\bibitem{ChenYu}
D.~Chen and F.~Yu,
An algebro-geometric perspective on the topology of moduli spaces of differentials,
arXiv:2601.13127.
\bibitem{Lanneau}
E.~Lanneau,
Parity of the spin structure defined by a quadratic differential,
arXiv:math/0210116.
\end{thebibliography}
\endgroup
\end{document}
